\clearpage
\fancyhead[L]{\small HARDING'S STAR ATLAS.}
\hypertarget{native-546-lower}{}\label{native-546-lower}
\begin{center}\small Göttingische gelehrte Anzeigen. No. 158. pp. 1569..1571. October 3, 1812.\\[3pt]\rule{.72\linewidth}{.3pt}\end{center}

The first two instalments of our Professor \textsc{Harding}'s great star atlas appeared in 1809 and 1810. Both were reviewed in our journal at the time. We now give the review of the third instalment, issued in 1811, together with the \SourceLeaf{547}review of the fourth, which has just appeared. We already described the plan and distinctive features of this work in detail when the first instalment appeared, so here we need only assure readers that the later instalments have been prepared in exactly the same spirit and with the same care, and in some respects even surpass the first in richness of content and fineness of engraving.

The \emph{third} instalment consists of sheets VII, XVI, XVII and XVIII. Sheet No. VII extends in right ascension from $239^\circ$ to $281^\circ$ and in declination from $2^\circ$ north to $32^\circ$ south. It thus contains parts of the constellations Scorpius Sagittarius, Ophiuchus, Serpens and Sobieski's Shield. Sheet No. XVI covers the same range of right ascension as No. VII and extends in declination from $1^\circ$ south to $33^\circ$ north. Besides most of Hercules, it contains parts of Serpens, Ophiuchus, Poniatovsky's Bull and Lyra. The exceptionally rich sheet No. XVII, for which Mr. \textsc{von Lindenau} supplied the author with a considerable number of stellar position determinations, adjoins the preceding sheet, has the same range of declination and extends to $321^\circ$ in right ascension. Besides all of Delphinus, Sagitta and the Fox with the Goose, it contains most of Equuleus and of the Eagle and Antinous, as well as parts of Poniatovsky's Bull, Serpens, Hercules, Lyra, Cygnus, Pegasus and Aquarius. Sheet No. XVIII in turn adjoins the preceding sheet in the same way and extends to a right ascension of $1^\circ$. On it we find nearly all of Pegasus, together with parts of Cygnus, Equuleus, Aquarius, Pisces and Andromeda.

The \emph{fourth instalment} gives us only southern zones, for which the author himself observed many additional stars during his stay at the Mannheim Observatory. It consists of sheets XI, XII, XIII and XIV. All these sheets extend from $1^\circ$ north to $33^\circ$ south and adjoin one another so as to cover the zone from $39^\circ$ to $201^\circ$ in right ascension. Thus they show, in part or in full, Cetus, George's Harp, the Laboratory, Eridanus, the Brandenburg Scepter, the Sculptor's Workshop, Orion, Lepus, Columba, Canis Major, the ship Argo, Monoceros, the Printer's Workshop, the Compass, the Air Pump, Hydra, Sextans, Leo, Virgo, Crater, Corvus and Centaurus. We need only add that the \textsc{Vandenhoeck} bookshop has now taken over the sale of the charts.
